%% Question bank chapter file - generated by bank_tools.py build (do not edit while the bank is still being built:
%% it is regenerated from the per-test data in Question Bank\data). Every question is one bankq block:
%%   \begin{bankq}[id=..., source={...}, paper=.., qno=.., type=..., subtopic={...}, answer={...}, answer_src=key|solved,
%%                 working={...}, difficulty=1-10, diff_parts={K S I C F}, diff_note={...},
%%                 group=<passage id or empty>, also_in={duplicate copies in other tests}]
%% difficulty = min(10, K+S+I+C+F-1): K concepts 1-3, S steps 1-4, I insight 0-2, C computation 0-2, F format 0-1
%%   ...question LaTeX, options, figures (\bfig / \bfigi -> img/)...
%%   \end{bankq}
%% Passages shared by several questions are bankpassage blocks; their questions carry group=<passage id>.
\documentclass[10pt]{article}
\usepackage[a4paper,margin=1.8cm]{geometry}
\usepackage[T1]{fontenc}
\usepackage{amsmath,amssymb,mathtools}
\usepackage{graphicx}
\usepackage[export]{adjustbox}
\usepackage{xcolor,needspace,keyval}
\usepackage{enumitem}
\usepackage{booktabs,array,tabularx}
\usepackage[hidelinks]{hyperref}
\setlength{\parindent}{0pt}
\setlength{\parskip}{3pt}
\makeatletter
\define@key{bq}{id}{\def\bq@id{#1}}\define@key{bq}{source}{\def\bq@src{#1}}\define@key{bq}{paper}{\def\bq@paper{#1}}
\define@key{bq}{qno}{\def\bq@qno{#1}}\define@key{bq}{type}{\def\bq@type{#1}}\define@key{bq}{subtopic}{\def\bq@sub{#1}}
\define@key{bq}{answer}{\def\bq@ans{#1}}\define@key{bq}{answer_src}{\def\bq@asrc{#1}}\define@key{bq}{working}{\def\bq@work{#1}}
\define@key{bq}{group}{\def\bq@group{#1}}\define@key{bq}{also_in}{\def\bq@also{#1}}
\define@key{bq}{difficulty}{\def\bq@diff{#1}}\define@key{bq}{diff_parts}{\def\bq@dparts{#1}}\define@key{bq}{diff_note}{\def\bq@dnote{#1}}
\newenvironment{bankq}[1][]{\def\bq@work{}\def\bq@also{}\def\bq@group{}\def\bq@diff{?}\def\bq@dparts{}\def\bq@dnote{}\setkeys{bq}{#1}%
  \par\needspace{4\baselineskip}\medskip\noindent{\small\textbf{\bq@id}\hfill\color{gray}[\bq@src, Paper \bq@paper, Q\bq@qno]
  \bq@type\ $\cdot$ \bq@sub\ $\cdot$ \textbf{D\bq@diff}}\par\nopagebreak\relax}
  {\par{\small\color{teal}Answer: \textbf{\bq@ans}\ifx\bq@work\@empty\else\ (solved: \bq@work)\fi
  \ifx\bq@also\@empty\else\par\color{gray}Also in: \bq@also\fi}\par}
\define@key{bp}{id}{\def\bp@id{#1}}\define@key{bp}{source}{\def\bp@src{#1}}\define@key{bp}{questions}{\def\bp@qs{#1}}
\newenvironment{bankpassage}[1][]{\setkeys{bp}{#1}\par\needspace{4\baselineskip}\bigskip\noindent
  {\small\textbf{Passage \bp@id}\hfill\color{gray}[\bp@src; questions \bp@qs]}\par\nopagebreak\relax\itshape}{\par}
\makeatother
\newcommand{\bfig}[2]{\par\smallskip{\centering\includegraphics[width=#1pt,max width=0.62\linewidth,max height=0.32\textheight]{img/#2}\par}\smallskip}
\newcommand{\bfigi}[2]{\includegraphics[width=#1pt,max width=\linewidth,max height=0.32\textheight,valign=c]{img/#2}}
\newcommand{\bmissing}[1]{\textit{[figure missing in the source conversion: \detokenize{#1}]}}
\newenvironment{qoptions}{\begin{itemize}[leftmargin=1.8em,itemsep=1pt,topsep=2pt,label={}]}{\end{itemize}}
\newcommand{\opt}[2]{\item[(#1)] #2}
\providecommand{\fwstart}{}\providecommand{\fwend}{}\providecommand{\qfullsplit}{}
\newcolumntype{L}[1]{>{\raggedright\arraybackslash}p{#1}}
\newcolumntype{C}[1]{>{\centering\arraybackslash}p{#1}}
\newcolumntype{R}[1]{>{\raggedleft\arraybackslash}p{#1}}
\providecommand{\sectionhdr}[1]{}\providecommand{\subsectionhdr}[1]{}


\begin{document}
{\Large\bfseries Intro to Vectors --- Question Bank}\par\smallskip
{\small Mathematics chapter M30. Questions from all converted test papers, grouped by source series. Answers: printed key where it matches the paper, otherwise solved (answer\_src=solved, with working). Exact repeats across papers appear once, with also\_in listing the other copies.}\par\medskip\hrule

\section*{Score-I papers}
\begin{bankq}[id=SI05-P1-MATH-S2-Q07, source={Score-I - Test 05 (22 Jan 2022)}, paper=1, qno=7, type=single, subtopic={Representation}, answer={B}, answer_src=key, difficulty=4, diff_parts={K2 S2 I1 C1 F0 adj-1}, diff_note={concepts: match range, a vector expression with integers | steps: ranges using unit vectors and triangle inequality; integer values in range; match}]
\par\nopagebreak[4]
\begin{center}\small\setlength{\tabcolsep}{3pt}
\begin{tabular}{|c|p{4.2cm}|c|c|}
\hline
\multicolumn{2}{|c|}{List-I} & \multicolumn{2}{c|}{List-II}\\ \hline
(P) & If $\vec a$ and $\vec b$ are any two unit vectors, then possible integer in the range of $\dfrac{3\left|\vec a+\vec b\right|}{2}+2\left|\vec a-\vec b\right|$ can be & (1) & $0$\\ \hline
(Q) & If $\hat a,\hat b$ and $\hat c$ are non coplanar vectors then the vectors $\vec V_1=\hat a+2\hat b+3\hat c$; $\vec V_2=\lambda\hat b+4\hat c$ and $\vec V_3=(2\lambda-1)\hat c$ where $\lambda$ is a scalar, can be coplanar, for $\lambda$ equal to & (2) & $1$\\ \hline
(R) & The distance of the z-axis from the image of the point A$(2,-3,3)$ in the plane $x-2y-z+1=0$, is & (3) & $2$\\ \hline
(S) & The figure shows a pyramid OABCD (where O is the origin). The base OABC is a square of side 1 unit. The point D is at a distance of 1 unit from the square base. D lies on the normal passing through the centroid of the base. If the angle between $\overrightarrow{\mathrm{OB}}$ and $\overrightarrow{\mathrm{OD}}$ is $\tan^{-1}\sqrt{K}$, then K is equal to\par\smallskip\bfigi{80.3}{SI05_p1_math_s2_q07_S.png} & (4) & $-1$\\ \hline
 & & (5) & $4$\\ \hline
\end{tabular}
\end{center}
\begin{qoptions}
\opt{A}{P $\to$ 4;Q $\to$ 2;R $\to$ 1;S $\to$ 5}
\opt{B}{P $\to$ 5;Q $\to$ 1;R $\to$ 2;S $\to$ 3}
\opt{C}{P $\to$ 5;Q $\to$ 1;R $\to$ 3;S $\to$ 2}
\opt{D}{P $\to$ 4;Q $\to$ 1;R $\to$ 2;S $\to$ 5}
\end{qoptions}
\end{bankq}
\begin{bankq}[id=SI05-P1-MATH-S3-Q07, source={Score-I - Test 05 (22 Jan 2022)}, paper=1, qno=7, type=numerical, subtopic={Position vectors}, answer={4.00}, answer_src=key, difficulty=5, diff_parts={K2 S2 I1 C1 F1 adj-1}, diff_note={concepts: vector identity for a quadrilateral, midpoints of diagonals | steps: sum of squares of sides = sum of squares of diagonals + 4 PQ^2; use given squared lengths; find asked}, also_in={TMAS - Test 04 (19 Nov 2023) P1 Q4}]
Let ABCD be any quadrilateral. P and Q are the midpoints of its diagonal. If $\left|\overrightarrow{AB}\right|^2+\left|\overrightarrow{BC}\right|^2+\left|\overrightarrow{CD}\right|^2+\left|\overrightarrow{DA}\right|^2-\left|\overrightarrow{AC}\right|^2-\left|\overrightarrow{BD}\right|^2=\lambda\left|\overrightarrow{PQ}\right|^2$, then find the value of $\lambda$.
\end{bankq}
\begin{bankq}[id=SI09-P2-MATH-S1-Q05, source={Score-I - Test 09 (16 Feb 2022)}, paper=2, qno=5, type=multi, subtopic={Position vectors}, answer={A,B,C}, answer_src=key, difficulty=5, diff_parts={K2 S2 I1 C1 F1 adj-1}, diff_note={concepts: vectors from origin, a plane or volume relation | steps: dot and cross products; area and volume formulas; check statements}]
Let $O$ be the origin and $\overrightarrow{OA}=2\hat i+2\hat j+\hat k$, $\overrightarrow{OB}=\hat i-2\hat j+2\hat k$ and $\overrightarrow{OC}=\dfrac12\left(\overrightarrow{OB}-\lambda\overrightarrow{OA}\right)$ for some $\lambda>0$. If $\left|\overrightarrow{OB}\times\overrightarrow{OC}\right|=\dfrac92$, then which of the following statements is (are) \textbf{TRUE} ?
\begin{qoptions}
\opt{A}{Projection of $\overrightarrow{OC}$ on $\overrightarrow{OA}$ is $-\dfrac32$}
\opt{B}{Area of the triangle $OAB$ is $\dfrac92$}
\opt{C}{Area of the triangle $ABC$ is $\dfrac92$}
\opt{D}{The acute angle between the diagonals of the parallelogram with adjacent sides $\overrightarrow{OA}$ and $\overrightarrow{OC}$ is $\dfrac\pi3$}
\end{qoptions}
\end{bankq}
\begin{bankq}[id=SI31-P2-MATH-S1-Q13, source={Score-I - Test 31 (29 Dec 2022)}, paper=2, qno=13, type=single, subtopic={Direction cosines}, answer={C}, answer_src=key, difficulty=5, diff_parts={K2 S2 I1 C1 F0}, diff_note={concepts: condition for two lines, given direction cosines to be parallel | steps: eliminate n from p l + q m + r n = 0; substitute in a l^2 + b m^2 + c n^2 = 0 to get homogeneous quadratic in l m; equal roots condition gives relation in p q r a b c}]
The direction cosines $\ell$, $m$, $n$ of two lines are given by $p\ell+qm+rn=0$ and $a\ell^2+bm^2+cn^2=0$ where $p$, $q$, $r$, $a$, $b$, $c$, are arbitrary constants. If these two lines are parallel then
\begin{qoptions}
\opt{A}{$\dfrac{b+c}{p^2}+\dfrac{c+a}{q^2}+\dfrac{a+b}{r^2}=0$}
\opt{B}{$\dfrac{a^2}{p}+\dfrac{b^2}{q}+\dfrac{c^2}{r}=0$}
\opt{C}{$\dfrac{p^2}{a}+\dfrac{q^2}{b}+\dfrac{r^2}{c}=0$}
\opt{D}{$p^2(b+c)+q^2(a+c)+r^2(a+b)=0$}
\end{qoptions}
\end{bankq}

\section*{Score-II papers}
\begin{bankq}[id=SII04-P1-MATH-S1-Q01, source={Score-II - Test 04 (07 May 2021)}, paper=1, qno=1, type=single, subtopic={Representation}, answer={C}, answer_src=key, difficulty=2, diff_parts={K1 S2 I0 C0 F0}, diff_note={concepts: statements about linear dependence of vectors | steps: non collinear vectors are independent; parallel vectors are dependent; four vectors in space dependent; choose the correct statement}, also_in={JEE Advanced (2021 Series) - Test 09 (07 May 2021) P1 Q1}]
Which of the following statements is correct:

\begin{qoptions}
\opt{A}{Two non collinear vectors are always linearly dependent}
\opt{B}{Two parallel non zero vectors are always linearly independent}
\opt{C}{Any four vectors in 3-d space are always linearly dependent}
\opt{D}{None of these}
\end{qoptions}
\end{bankq}
\begin{bankq}[id=SII15-P2-MATH-S3-Q05, source={Score-II - Test 15 (09 Feb 2024)}, paper=2, qno=5, type=numerical, subtopic={Direction cosines}, answer={1}, answer_src=key, difficulty=5, diff_parts={K1 S2 I1 C1 F1}, diff_note={concepts: slope angle of a line to meet x = -7 with a distance condition | steps: parametric line from P; distance to the point on x = -7; solve for tan theta}]
A line which makes an acute angle `$\theta$' with positive direction of x-axis is drawn from $P(5,4)$ to meet the line $x=-7$ at $Q$ and $y=9$ at $R$ then $\dfrac{144}{PQ^2}+\dfrac{25}{PR^2}$ is
\end{bankq}
\begin{bankq}[id=SII17-P1-MATH-S1-Q16, source={Score-II - Test 17 (13 Feb 2024)}, paper=1, qno=16, type=single, subtopic={Direction cosines}, answer={C}, answer_src=solved, working={Direction cosines (1/sqrt2,1/2,1/2); segment vector (0,2,-3): projection 1-3/2=-1/2, magnitude 1/2.}, difficulty=4, diff_parts={K1 S2 I1 C1 F0}, diff_note={concepts: projection of a segment on a line with given direction angles | steps: direction cosines from the angles; vector between points; dot product with the unit direction}]
A line makes angle $45^\circ$, $60^\circ$, and $60^\circ$ with the coordinate axis, the projection of line segments on line which joins point ($-$ 1,2,3) and ($-$ 1, 4, 0) are-
\begin{qoptions}
\opt{A}{3/2}
\opt{B}{1/3}
\opt{C}{1/2}
\opt{D}{2/3}
\end{qoptions}
\end{bankq}
\begin{bankq}[id=SII19-P1-MATH-S1-Q20, source={Score-II - Test 19 (20 Feb 2024)}, paper=1, qno=20, type=single, subtopic={Direction cosines}, answer={A}, answer_src=key, difficulty=4, diff_parts={K1 S2 I1 C1 F0}, diff_note={concepts: perpendicular direction cosines condition | steps: eliminate n from the two relations; quadratic in l/m; product of roots gives the perpendicularity condition; compare options}]
The straight lines whose direction cosines are given by $a\ell+bm+cn=0$, $fmn+gn\ell+h\ell m=0$ are perpendicular, if :-
\begin{qoptions}
\opt{A}{$\dfrac fa+\dfrac gb+\dfrac hc=0$}
\opt{B}{$\sqrt{\dfrac af}+\sqrt{\dfrac bg}+\sqrt{\dfrac ch}=0$}
\opt{C}{$\sqrt{af}=\sqrt{bg}=\sqrt{ch}$}
\opt{D}{$\sqrt{\dfrac af}=\sqrt{\dfrac bg}=\sqrt{\dfrac ch}$}
\end{qoptions}
\end{bankq}
\begin{bankq}[id=SII29-P1-MATH-S3-Q10, source={Score-II - Test 29 (26 Mar 2024)}, paper=1, qno=10, type=single, subtopic={Position vectors}, answer={C}, answer_src=key, difficulty=5, diff_parts={K2 S2 I1 C1 F0}, diff_note={concepts: ratio theorems on cevians, trigonometric equation with vectors | steps: P from mass point or Menelaus with the given ratios; Q solve sin 2 theta over cos theta equal to the ratio and count roots in 0 to 4 pi; R from the vector relation using the BD and BE ratios; S from the same ratio}]
In the adjoining figure

\begin{center}
\bfigi{135.4}{SII29_p3math_q10_triangle.png}
\end{center}

BD : DC = 3 : 4,
BP : PE = $\sin2\theta$ : $\cos\theta$ \& E is the mid point of AC.

\textbf{List-I}\\
(P) Value of $\left[\dfrac{AP}{PD}\right]$ is (where [.] denotes greatest integer function)\\
(Q) Number of values of $\theta\in[0,4\pi]$ is\\
(R) If $\lambda(\overrightarrow{BA}+\overrightarrow{BC})=10\overrightarrow{BP}$, then $\lambda$ is\\
(S) Value of $\left[\dfrac{BP}{PE}\right]$ is (where [.] denotes greatest integer function)

\textbf{List-II}\\
(1) 3\quad (2) 2\quad (3) 1\quad (4) 4\quad (5) 6

\begin{qoptions}
\opt{A}{P $\to$ 3; Q $\to$ 1; R $\to$ 5; S $\to$ 2}
\opt{B}{P $\to$ 1; Q $\to$ 2; R $\to$ 3; S $\to$ 5}
\opt{C}{P $\to$ 2; Q $\to$ 4; R $\to$ 1; S $\to$ 3}
\opt{D}{P $\to$ 4; Q $\to$ 3; R $\to$ 5; S $\to$ 4}
\end{qoptions}
\end{bankq}
\begin{bankq}[id=SII31-P1-MATH-S1-Q13, source={Score-II - Test 31 (29 Mar 2024)}, paper=1, qno=13, type=single, subtopic={Direction cosines}, answer={C}, answer_src=key, difficulty=5, diff_parts={K1 S3 I1 C1 F0}, diff_note={concepts: angle between lines from direction cosine equations | steps: n = 2l + 2m; substitute into mn + nl + lm; factor; two direction sets; angle}]
The angle between the straight lines, whose direction cosines $\ell$, $m$, $n$ are given by the equations $2\ell+2m-n=0$ and $mn+n\ell+\ell m=0$ is $\dfrac\pi k$, then the sum of all the divisors of $6k^2$
\begin{qoptions}
\opt{A}{120}
\opt{B}{30}
\opt{C}{60}
\opt{D}{90}
\end{qoptions}
\end{bankq}

\section*{Score Advanced papers}
\begin{bankq}[id=SA01-P2-MATH-S1-Q07, source={Score Advanced - Test 01 (22 Sep 2020)}, paper=2, qno=7, type=multi, subtopic={Position vectors}, answer={B,C,D}, answer_src=key, difficulty=5, diff_parts={K1 S2 I1 C1 F1}, diff_note={concepts: centroid properties and distance sum | steps: PA^2+PB^2+PC^2 identity with G; minimum at G; check}]
Let $G$ be the centroid of the $\Delta ABC$ whose sides are of length $a$, $b$, $c$. If $P$ be a point in the plane of $\Delta ABC$ such that $PA=1$, $PB=3$, $PC=4$ and $PG=2$, then the value of $a^2+b^2+c^2$ is less than :
\begin{qoptions}
\opt{A}{40}
\opt{B}{44}
\opt{C}{46}
\opt{D}{48}
\end{qoptions}
\end{bankq}
\begin{bankq}[id=SA09-P2-MATH-S1-Q08, source={Score Advanced - Test 09 (29 Apr 2023)}, paper=2, qno=8, type=multi, subtopic={Position vectors}, answer={C,D}, answer_src=key, difficulty=6, diff_parts={K1 S3 I1 C1 F1}, diff_note={concepts: collinearity condition in a parallelogram  a transversal | steps: AQ = AP + AR (parallelogram diagonal); points P Q R collinear gives 1 = 1/(alpha ) relation; AQ = gamma AC; relation among alpha beta gamma; check statements}]
Consider a parallelogram ABCD. A straight line is drawn which cuts AB, AC, AD at points P, Q, R respectively. If $\overrightarrow{\text{AP}}$ = $\alpha\overrightarrow{\text{AB}}$, $\overrightarrow{\text{AR}}$ = $\beta\overrightarrow{\text{AD}}$ and $\overrightarrow{\text{AQ}}$ = $\dfrac{\gamma}{2}\overrightarrow{\text{AC}}$ where $\alpha$, $\beta$, $\gamma$ are positive scalars, then
\begin{qoptions}
\opt{A}{$\alpha$, $\beta$, $\gamma$ are in A.P.}
\opt{B}{$\alpha$, $\beta$, $\gamma$ are in H.P.}
\opt{C}{$\alpha$, $\gamma$, $\beta$ are in H.P.}
\opt{D}{$\alpha\beta\geq\gamma^2$}
\end{qoptions}
\end{bankq}
\begin{bankq}[id=SA12-P1-MATH-S3-Q03, source={Score Advanced - Test 12 (11 May 2023)}, paper=1, qno=3, type=numerical, subtopic={Representation}, answer={4}, answer_src=key, difficulty=5, diff_parts={K1 S2 I1 C1 F1}, diff_note={concepts: vector rotation and projection | steps: rotate V1 about an axis; coordinates; dot or cross; compute}, also_in={Score Advanced - Test 29 (01 May 2024) P2 Q3}]
Let a vector $\bar{\text{V}}_1$ = $\overline{\text{OP}}$, where O is origin and P is (2, 0,2). Vector $\bar{\text{V}}_1$ is rotated about origin by an angle of $\dfrac{\pi}{3}$ in a plane which is perpendicular to XOZ plane to get vector $\bar{\text{V}}_2$ where $\bar{\text{V}}_2$. $\hat{\text{j}}$ $>$ 0. Again $\bar{\text{V}}_1$ is rotated about origin by right angle in XOZ plane to get vector $\bar{\text{V}}_3$. If V is volume of tetrahedron whose coterminous edges are $\bar{\text{V}}_1$, $\bar{\text{V}}_2$ and $\bar{\text{V}}_3$ then value of $\dfrac{3\text{V}}{\sqrt{6}}$ is.
\end{bankq}
\begin{bankq}[id=SA14-P1-MATH-S1-Q05, source={Score Advanced - Test 14 (16 May 2023)}, paper=1, qno=5, type=multi, subtopic={Position vectors}, answer={A,B,C,D}, answer_src=key, difficulty=5, diff_parts={K1 S2 I1 C1 F1}, diff_note={concepts: collinearity  position of points from c = x a + y b | steps: affine combination x + y = 1 means C on line AB; sign conditions determine segment or extensions; x + y with other values; check four}, also_in={Score Advanced - Test 20 (06 Apr 2024) P1 Q5}]
Let position vectors of points A, B, C are $\vec{\text{a}},\vec{\text{b}},\vec{\text{c}}$ respectively where $\vec{\text{c}}$ = x$\vec{\text{a}}$ + y$\vec{\text{b}}$
\begin{qoptions}
\opt{A}{If x + y = 1, x $>$ 0, y $<$ 0, then C lies on line AB but not on line segment AB}
\opt{B}{If x $>$ 0, y $>$ 0, x + y $<$ 1, then C lies inside the triangle $\Delta$ABO (O is origin)}
\opt{C}{If x $<$ 0, y $>$ 0 and x + y $>$ 1, then on plane containing points O, A, B, the points C and origin (O) lies on opposite side of line AB.}
\opt{D}{If x $\leq$ 0, y $\leq$ 0 and x + y = $-$ 1, then locus of point C is reflection of line segment AB about origin.}
\end{qoptions}
\end{bankq}
\begin{bankq}[id=SA16-P1-MATH-S4-Q09, source={Score Advanced - Test 16 (20 May 2023)}, paper=1, qno=9, type=numerical, subtopic={Position vectors}, answer={4}, answer_src=key, difficulty=6, diff_parts={K1 S3 I1 C1 F1}, diff_note={concepts: point on the line joining face centres  a parallelepiped with a vector condition | steps: M and M' centres of opposite faces; position vectors; vector sum condition sets O as a weighted point on MM'; ratio MO:OM'; evaluate the integer asked}]
If ABCDEFGH is a parallelopiped and $M$ and $M'$ are the point of intersections of diagonals of parallelograms ABCD and EFGH. If $O$ is a point on $MM'$ such that $\overrightarrow{OA}+\overrightarrow{OB}+\overrightarrow{OC}+\overrightarrow{OD}=\overrightarrow{M'O}$ and $\lambda\overrightarrow{MO}=\overrightarrow{OM'}$, then $\lambda$ is \underline{\hspace{1.5cm}}.
\end{bankq}
\begin{bankq}[id=SA24-P1-MATH-S1-Q06, source={Score Advanced - Test 24 (17 Apr 2024)}, paper=1, qno=6, type=multi, subtopic={Representation}, answer={A,B}, answer_src=key, difficulty=5, diff_parts={K1 S2 I1 C1 F1}, diff_note={concepts: range  the half angle given distance between unit vectors | steps: abs(a - b) = 2 sin(theta) < 1 gives sin theta < 1/2; theta in [0,pi]; solution sets [0,pi/6) and (5pi/6, pi]}]
If the unit vectors $\vec a$ and $\vec b$ are inclined at an angle $2\theta$ such that $|\vec a-\vec b|<1$ and $0\leqslant\theta\leqslant\pi$ then $\theta$ lies in the interval
\begin{qoptions}
\opt{A}{$[0,\pi/6)$}
\opt{B}{$(5\pi/6,\pi]$}
\opt{C}{$(\pi/6,\pi/2]$}
\opt{D}{$[\pi/2,5\pi/6)$}
\end{qoptions}
\end{bankq}
\begin{bankq}[id=SA30-P1-MATH-S1-Q04, source={Score Advanced - Test 30 (04 May 2024)}, paper=1, qno=4, type=single, subtopic={Representation}, answer={D}, answer_src=key, difficulty=5, diff_parts={K1 S3 I1 C1 F0}, diff_note={concepts: angle between vectors depending on a parameter  projection condition | steps: v1 = (a^2, -2, a+1), v2 = (2, ); condition for angle acute/obtuse or coplanarity; scalar product as a quadratic in a; solve inequality; count integer values}, also_in={Score Advanced - Test 31 (04 May 2024) P1 Q4}]
Let $\vec{\mathrm v}_1$ = a$^2\hat{\mathrm i}$ $-$ 2$\hat{\mathrm j}$ + (a + 1)$\hat{\mathrm k}$ and $\vec{\mathrm v}_2$ = 2$\hat{\mathrm i}$ + $\hat{\mathrm j}$ $-$ (a$^2$ $-$ 1)$\hat{\mathrm k}$ be two vectors such that $\vec{\mathrm v}_2$ lies in the plane 3x $-$ $\lambda$y + z = 1 and ($\vec{\mathrm v}_1$ + $\vec{\mathrm v}_2$).$\vec{\mathrm v}_1$ = a$^4$ + a$^2$ + 2a + 5 then greatest volume of the parallelopiped whose co-terminous edges are represented by $\vec{\mathrm v}_1$, $\vec{\mathrm v}_2$ and $\vec{\mathrm n}$ where $\vec{\mathrm n}$ is the normal vector of the given plane, is -
\begin{qoptions}
\opt{A}{3}
\opt{B}{5}
\opt{C}{15}
\opt{D}{25}
\end{qoptions}
\end{bankq}

\section*{Pre-Score (2023-24) papers}
\begin{bankq}[id=PS01-P2-MATH-S4-Q07, source={Pre-Score (2023-24) - Test 01 (08 Oct 2023)}, paper=2, qno=7, type=numerical, subtopic={Position vectors}, answer={7}, answer_src=key, working={Printed key.}, difficulty=6, diff_parts={K1 S2 I1 C2 F1}, diff_note={concepts: orthocentre of medial triangle and distance | steps: the medial triangle is similar with orthocentre equal to the circumcentre of the original triangle; compute circumcentre of ABC in space; distance to Q | I: the orthocentre of the medial triangle is the circumcentre of ABC | C: finding the circumcentre in 3D coordinates}]
Let position vectors of vertices of $\Delta ABC$ be $3\hat i+4\hat j+5\hat k$, $-2\hat i+6\hat j+\sqrt{10}\hat k$, $5\hat j-5\hat k$. Let $DEF$ be triangle formed by joining the mid points of sides of $\Delta ABC$. If $\vec P$ is the position vector of orthocentre of $\Delta DEF$ and position vector of $Q$ is $2\hat i-3\hat j+6\hat k$, then find $\left|\overrightarrow{PQ}\right|$.
\end{bankq}
\begin{bankq}[id=PS01-P2-MATH-S4-Q09, source={Pre-Score (2023-24) - Test 01 (08 Oct 2023)}, paper=2, qno=9, type=numerical, subtopic={Position vectors}, answer={2}, answer_src=key, working={Printed key.}, difficulty=5, diff_parts={K1 S2 I1 C1 F1}, diff_note={concepts: collinearity of combined vectors | steps: the two combinations are collinear so coefficients are proportional; compare coefficients of a b c in r1 plus 2 r2 and 2 r1 plus r2; solve for p and q; compute p squared plus q squared | I: non coplanar vectors give three equal ratio equations | C: solving the resulting equations}]
Let $\vec a,\vec b,\vec c$ are three non-zero, non-collinear vectors and $\vec r_1=p\vec a+q\vec b+\vec c$ and $\vec r_2=\vec a+p\vec b+q\vec c$. If the vectors $\vec r_1+2\vec r_2$ and $2\vec r_1+\vec r_2$ are collinear then find value of $p^2+q^2$.
\end{bankq}

\section*{JEE Advanced (2021 Series) papers}
\begin{bankq}[id=JA05-P1-MATH-S3-Q01, source={JEE Advanced (2021 Series) - Test 05 (11 Apr 2021)}, paper=1, qno=1, type=numerical, subtopic={Representation}, answer={0.00}, answer_src=key, difficulty=5, diff_parts={K1 S2 I2 C0 F1}, diff_note={concepts: linear independence of three vectors | steps: x, y and x cross y are linearly independent so all three coefficients vanish; a quadratic with three distinct arguments alpha beta gamma has three roots so a = b = c = 0; sum 0 | I: independence of x y and their cross product forces each coefficient to zero}]
Let $\vec x,\vec y$ are non zero and non collinear vectors such that
\[(a\alpha^2+b\alpha+c)\vec x+(a\beta^2+b\beta+c)\vec y+(a\gamma^2+b\gamma+c)(\vec x\times\vec y)=0,\]
then $(a+b+c)$ is (where $a,b,c\in\mathbb R$)
\end{bankq}
\begin{bankq}[id=JA14-P1-MATH-S1-Q05, source={JEE Advanced (2021 Series) - Test 14 (06 Jun 2021)}, paper=1, qno=5, type=single, subtopic={Position vectors}, answer={A}, answer_src=solved, working={Collinearity of B1, D1 and C1 gives 1/lambda3 = 1/lambda1 + 1/lambda2.}, difficulty=3, diff_parts={K1 S2 I1 C0 F0}, diff_note={concepts: collinearity in a parallelogram with vectors | steps: write position vectors of B1, C1, D1 in terms of AB and AD; collinearity gives a linear relation; compare coefficients to get reciprocals | I: collinearity of three points on line L is expressed through the relation among reciprocals}]
A line $L$ cuts the lines $AB$, $AC$ and $AD$ of a parallelogram $ABCD$ at points $B_1$, $C_1$ and $D_1$ respectively. If $\overrightarrow{AB_1}=\lambda_1\overrightarrow{AB}$, $\overrightarrow{AD_1}=\lambda_2\overrightarrow{AD}$, $\overrightarrow{AC_1}=\lambda_3\overrightarrow{AC}$ then $\dfrac1{\lambda_3}=$
\begin{qoptions}
\opt{A}{$\dfrac1{\lambda_1}+\dfrac1{\lambda_2}$}
\opt{B}{$\dfrac1{\lambda_1}-\dfrac1{\lambda_2}$}
\opt{C}{$-\lambda_1+\lambda_2$}
\opt{D}{$\lambda_1+\lambda_2$}
\end{qoptions}
\end{bankq}

\section*{TOAS papers}
\begin{bankq}[id=TOAS17-P2-MATH-S2-Q06, source={TOAS - Test 17 (18 Sep 2022)}, paper=2, qno=6, type=numerical, subtopic={Direction cosines}, answer={9}, answer_src=solved, working={Three unit vectors with equal pairwise angles are coplanar at the maximum, $\theta=2\pi/3=\pi/k$, so $k=3/2$ and $4k^2=9$}, difficulty=4, diff_parts={K1 S2 I1 C0 F1}, diff_note={concepts: angles between unit vectors | steps: equal pairwise angles; maximum is coplanar; k}]
If $|\vec a|=|\vec b|=|\vec c|=1$ and $\vec a\mathbin{\wedge}\vec b=\vec c\mathbin{\wedge}\vec b=\vec a\mathbin{\wedge}\vec c=\theta$ where $\vec p\mathbin{\wedge}\vec q$ represents the angle between vectors $\vec p$ and $\vec q$ then maximum value of $\theta$ is $\dfrac\pi k$. The value of $4k^2$ is
\end{bankq}
\begin{bankq}[id=TOAS17-P2-MATH-S2-Q10, source={TOAS - Test 17 (18 Sep 2022)}, paper=2, qno=10, type=numerical, subtopic={Representation}, answer={1}, answer_src=solved, working={Comparing coefficients: $2x-y=5$, $x-2y=4$ gives $x=2$, $y=-1$, so $x+y=1$}, difficulty=2, diff_parts={K1 S1 I0 C0 F1}, diff_note={concepts: equality of vectors | steps: compare coefficients; solve two equations}]
If $\vec a$ and $\vec b$ are non collinear vectors such that $(2x-y)\vec a+4\vec b=5\vec a+(x-2y)\vec b$ then $(x+y)$ is
\end{bankq}
\end{document}
